\documentclass[11pt, a4paper]{article}
\usepackage[nochapters]{classicthesis}                              % template
\usepackage[margin=42mm]{geometry}                                  % margins
\usepackage[utf8]{inputenc}                                         % allow utf-8 input
\usepackage[T1]{fontenc}                                            % use 8-bit T1 fonts
\usepackage{graphicx}                                               % images
\usepackage{url}                                                    % URL typesetting
\usepackage{booktabs}                                               % good-looking tables
\usepackage{multirow}                                               % for tables
\usepackage{amsfonts}                                               % blackboard math symbols
\usepackage{amsmath}                                                % math ops
\usepackage{nicefrac}                                               % compact 1/2, etc.
\usepackage{microtype}                                              % microtypography
\usepackage{float} %to force figures to appear at a specific location

\definecolor{darkblue}{rgb}{0, 0, 0.5}                              % define link color
\hypersetup{colorlinks=true,citecolor=darkblue,                     % set link color
            linkcolor=darkblue, urlcolor=darkblue}

% Your packages here
% \usepackage{...}                                                  % some info, maybe

% !!! PLEASE CHANGE THESE VARIABLES TO MATCH YOUR INFORMATION !!!

\def\thesistitle{From Household Decisions to Electric Grid Patterns}
                      % title
\def\subtitle{An Agent-Based Analysis of Consumer Heterogeneity and
Load Synchronization in Variable-Priced Electricity Systems}             % subtitle
    % ^if there is no subtitle, replace by \def\subtitle{}  
\def\yourname{Andrey Geschiere}                                                                       % ^first and last nam
 \def\yourprogramme{Cognitive Science \& Artificial Intelligence}    % uncomment this
\def\yourstudentnumber{2120505}                                      % ANR (or u-number), -> Anderen doen SNR
\def\finalmonth{January}
\def\finalyear{2000}
\def\supervisor{dr. Travis Wiltshire}                           % check the title(s) of your supervisor
\def\committee{dr. The Second Reader}                               % check the title(s) of your second reader
\def\acknowledgments{I want to express my gratitude towards my supervisor, Dr. Travis J. Wiltshire, for asking the right questions and steering this thesis towards completion with needed advice. I would also like to thank my thesis group members, Bram de Waal and Soham Desriaux, for providing an outside perspective and for recognizing progress that I often overlooked while immersed in the details of the project. Finally, I want to thank my parents, roommates, and Tilburg University Library personnel for providing the perfect environment to focus on this project.}
\def\wordcount{max 8,000 words, excluding references and appendices, ADD HERE}
% METADATA

\hypersetup{pdfauthor   = {Andrey Geschiere},
            pdftitle    = \thesistitle\ \subtitle,      %TEDOEN
            pdfsubject  = \CognitiveScienceandArtificialIntelligence\ Bachelor Thesis}

% CHOOSE EITHER OF -----------------------------------------
% IEEE STYLE:

% \usepackage[square,numbers]{natbib}                               % bracket-style refs
% \usepackage{natbib}                                               % OR: parenthesized refs
% \bibliographystyle{IEEEtranN}

% OR -------------------

\usepackage[natbibapa]{apacite}                                     % only parentheses
\bibliographystyle{apacite}                                         % 'cause APA

% !!! ------------------------------------------------------ !!!

\begin{document}
\input{frontmatter.tex}  % don't remove this :)

% --- start writing below:

\begin{abstract}
The concentration of residential demand into morning and evening peaks creates a temporal mismatch with the generation of renewable energy. This challenges grid operators, who have to integrate green power without compromising stability. As variable-priced contracts are spreading across Western Europe, it is necessary for effective policymaking to understand how heterogeneous household behaviors translate into aggregate demand patterns under these new pricing structures. This thesis examines these dynamics through an Agent-Based Model (ABM) of heterogeneous households, simulated over a 30-day period at 15-minute granularity. Household-level energy usage is modeled as collections of appliances with distinct power consumption and flexibility characteristics. Three behavioral architectures influence how agents schedule these appliances: habit-driven agents who concentrate usage on preferred times, price-responsive agents who shift appliance usage toward hourly price minima, and socially-influenced agents who gradually adopt observed schedules of their local network. The model is calibrated against Dutch grid data from the OpenSTEF Liander 2024 dataset. Results indicate that behavioral architecture is the dominant driver of individual-level outcomes within the model, with price-responsive agents shifting approximately four times more per day than habit-driven agents and achieving lower energy costs. At the system level, population composition explained over 83\% of rank variance across all tested metrics. Socially-influenced agents required sufficient exposure to price-responsive peers to shift meaningfully, pointing to a behavioral tipping dynamic that is relevant for demand response design.
\end{abstract}

\section{Data Source, Ethics, Code, and Technology statement}
Data from the Liander OpenSTEF 2024 dataset was used to build model components and calibration. Additional data from Yilmaz et al. (2017), Robinson et al. (2013) and Williams et al. (2025) was used to define empirical baselines in the code. These sources are cited where necessary in both the code and this thesis report, for further detail on data usage refer to \ref{subs:data}. The project code and analysis notebooks can be found  in the GitHub repository (ADD SOURCE TO GIT like Geschiere, 2026). No consent was required to use any of the data in this project. Analysis is performed on synthetic data, which was generated with the model implemented in this thesis. Simulation parameters are derived from cited research, and all figures and tables were created by the author. For a full package and version overview, refer to \ref{subs:technical} and the requirements.txt file in the repository. The thesis follows the LaTeX template provided by Tilburg University. Reference management was done through Zotero (Corporation for Digital Scholarship, 2025). Generative AI tools (Gemini (Google, 2026), Claude (Anthropic, 2026), ChatGPT (OpenAI, 2026)) were used for bugfixing and paraphrasing. Claude wrote pseudocode for the House (2014) network implementation as described in \ref{subs:networks}. Grammarly was used for grammar and spelling checks. 

\section{Introduction} \label{sec:introduction}

As global electricity use increases and the integration of sustainable energy becomes more urgent, maintaining the stability and reliability of electricity grids becomes an increasingly complex challenge (Akdemir et al., 2025; Dora et al., 2025; Siebert et al., 2017). This complexity is strengthened by the misalignment between the generation of renewable energy sources (RES) and concentrated consumer demand, resulting in damaging load peaks, avoidable maintenance costs, and local grid shutdowns (Eid et al., 2016; Tostado-Véliz et al., 2021). Such demand peaks are a result of synchronization, which is defined as the temporal and spatial alignment of behavior among interacting components within a complex system (Hudson et al., 2022). These synchronous patterns emerge within the electricity context as individual behaviors align across predictable timeframes. First, a morning ramp-up occurs as households awake and start their daily routines. A demand minimum is then observed during working hours, reflecting the collective shift away from the residential setting. The cycle concludes with a global maximum in the evening, caused by the collective activation of high-power appliances and electric vehicle chargers as people return to their homes. 

This rhythm in residential demand is predictable and generalizes across Western European countries (Anvari et al., 2022; Liander, 2024). Despite this predictability, the temporal disconnect between surplus RES and aggregate demand poses a challenge for Distribution System Operators (DSOs) to achieve an optimized, cost-effective, and future-proof usage of excess renewable power. Consequently, nonrenewable fuels still constitute the majority of the European energy supply (Eurostat, 2026). Given the European Union’s (EU) legal commitment to achieve climate neutrality by 2050 and its identification of "energy efficiency" as a core pillar of this transition, fulfilling this target requires addressing the behavior that drives synchronization (Consilium, 2026). Ultimately, it is these behavioral factors that compound into this current state of "energy inefficiency". The EU recognized this with their "Energy Efficiency First Paradigm", which promotes the opportunities of demand response (European Climate Foundation, 2016). Demand Response (DR), as defined by Wohlfarth et al. (2020), refers to the changing of consumption patterns to optimise the usage of energy supply. With DR playing a key role in the RES transition, many countries employ various strategies to change consumer behavior. 

A fundamental prerequisite to these strategies is the widespread installation of smart meters. A smart meter records energy usage data at a given interval and provides usage data to both consumers and DSOs (Zhou \& Brown, 2017). Installation of these devices facilitates the collection of the granular data necessary for variable pricing. In numerous European countries, such as Sweden, Norway, Denmark, Germany, and the Netherlands, widespread smart meter roll-outs pave the way for a shift towards variable priced energy systems (Eid 2016; Zhou \& Brown, 2017; Rijksinspectie Digitale Infrastructuur, 2025). Variable energy contract architectures can differ vastly, ranging from dynamic or static hour-by-hour tracking to seasonal price adjustments. Regardless of the specific structure, these contracts generally create a win-win scenario for both parties (Johanning et al., 2024). The consumer lowers their bill by shifting consumption to periods of high renewable generation, allowing the DSO to utilize surplus energy. Consequently, the consumer draws less power during the evening peak, which relieves the grid operator of meeting a demand surplus when the system is at its most strained. Despite their growing relevance, the share of households on a dynamic contract form a market minority, showing a 27\% adoption rate across Europe (ACER-CEER, 2024).  However, with these promising benefits and institutional support, adoption of dynamic pricing constructs is on an upward trajectory (Johanning et al., 2024). However, the aggregate grid effect of this transition is likely not be homogeneous, as is likely shaped by the composition of the adopting population and the diversity of behavioral responses within it. A grid populated predominantly by habitual consumers might respond very differently to dynamic pricing than one where a population is more aware of DR benefits.

Understanding the implications of such a transition, while keeping the heterogeneous behavioral drivers in mind, requires a methodology that can capture both system- and household-level behavior. Agent Based Modeling (ABM) provides an ideal framework for analyzing both scopes within a complex system (Wilensky \& Rand, 2015). The feasibility of this methodology within the energy context is demonstrated by Williams et al. (2025), who demonstrated that neighborhood-level electricity demand can be modeled with high fidelity, often staying within a 1-4\% margin of actual data. Plenty of case studies investigate ABM feasibility in simulating and predicting demand, but research into the effects of heterogeneous- and behavioral aspects remain limited.

\medskip

This study will examine the effects of diverse population compositions and the behavioral roles within those populations. Specifically distinguishing between three interacting behavioral profiles: 

(1) Habit-driven agents, who concentrate their usage on times most familiar to them with, little DR opportunity. 

(2) Price-responsive agents, who actively shift their load towards times of local pricing minima. 

(3) Socially-influenced agents, less strict habitual agents with the ability to gradually shift towards adaptation of price-responsive behavior, if this is observed within their network.  

\medskip

This agent-based framework captures the dynamics between ingrained habits and external DR influences aimed at behavioral change. It provides the means to see how heterogeneous households aggregate into system-wide effects. To analyze these dynamics, this thesis asks the following research questions:

\begin{itemize}
    \item[RQ1] \emph{How do different household behavioral decision
architectures (habit-driven, price-responsive,
socially-influenced) affect flexibility, adjustment, and cost-comfort
trade-offs within an ABM?}
    \item[RQ2] \emph{How do different configurations of heterogeneous households influence aggregate load, flexibility patterns, costs, and overall grid stability at the system level?}
\end{itemize}

By answering these research questions, this thesis aims to bridge a scientific gap between technical grid modeling and behavioral science. Societally, these findings provide insights for DSOs and policymakers as Europe necessarily transitions towards dynamic pricing. Ultimately, this research could inform the design of demand response strategies for establishing a future-proof electricity grid.

\section{Related Work}

\subsection{The Paradigm Shift} \label{subs:shift}
The architecture of electricity systems has historically followed the assumption that supply must be continuously adjusted to match consumption. Eid et al. (2016) formalize this as the \textit{generation-follows-demand} perspective. They define it as a structure where "generation units are synchronously operated to automatically supply the actual demand of electricity". Historically, this meant that (fossil) power generators bore the burden of maintaining grid stability by adjusting their output in real-time.

The penetration of weather-dependent RES within the electric grid has inverted this logic onto the consumer. Wind and solar generation are stochastic, producing electricity when conditions allow rather than when demand requires it. Eid et al. (2016) follow up by identifying the structural consequence, noting that "in future electricity systems the generation-follows-demand perspective is increasingly replaced by a demand-follows-generation perspective".  

A central feature of this new logic is the evolution of users from passive consumers into \textit{prosumers} who act as active participants within the demand chain. Parag and Sovacool (2016) describe prosumers as agents that "actively manage their own consumption and production of energy," relying on solar panels and smart-meter infrastructure to participate on both sides of the energy market. The growth of such distributed generation introduces variable local surpluses and deficits at the distribution level that conventional grid management is not designed to absorb.

According to Eid et al. (2016), dynamic pricing serves as the commercial instrument through which this paradigm shift reaches individual households. When wholesale prices vary hourly, households enrolled in variable contracts have a direct financial incentive to shift deferrable consumption toward periods of high renewable generation and low price. This direction has been gaining institutionalization across Europe. Johanning et al. (2024) note that the German GNDEW law now mandates that all energy providers must offer contracts with dynamic energy tariffs in their portfolio. This marks a significant mainstreaming of the demand-follows-generation logic through national policy. 

\subsection{Defining Usage Behavior} \label{subs:userbehav}
The demand response literature establishes that households respond to price signals and social nudges. However, this response is not uniform and not for the same reasons. Johanning et al. (2024) identify this as a structural point in the energy transition, noting that "households are not homogeneous in their electricity demand and consumption habits" and that "differences in socio-economic or attitudinal factors, as well as household composition, can influence electricity consumption and attitudes towards shifting electricity usage." These mixing components result in a heterogeneous population with each household having an unique stance against DR. The separation of these behavioral drivers by Johanning et al. (2024) motivated this thesis' approach for settling on the three behavioral types that define a behavioral profile. The selection of habitual, price-sensitive, and socially-influenced agents was motivated by the literature discussed in this subsection.

Habit persistence is the first and most prevalent mechanism. Häseler and Wulf (2024) regard it as the central obstacle to price-based DR. Even when smart meters and dynamic tariffs are available, consumers have been rooted in long-standing routines for generations. These behavioral patterns have been established in times of fixed electricity prices being the norm. Consequently, there is an absence of a societal norm or 'culture' regarding a positive scheduling of daily electrical appliance use. This behavioral observation defines the habit-driven agent in this model. Habitual agents concentrate consumption on times familiar to them, with little responsiveness to price signals.

Price responsiveness functions as the behavioral counterpart: the agent who actively shifts consumption toward cheaper periods in response to market signals. Häseler and Wulf (2024) characterize this profile, mentioning that "a household's financial savings potential depends on its responsiveness: the size of the loads to be shifted, by how long they can be shifted, and the consumer's proficiency at scheduling consumption based on the price profile."

Social influence constitutes the third mechanism in this ABM. Johanning et al. (2024) explicitly identify social network effects as an important modeling consideration in ABM energy demand. Empirical evidence supports the weight of this channel. Du et al. (2024) examined the influence of social networks on the adoption of energy-efficient replacements for inefficient appliances. They found that approximately 80\% of respondents in low-barrier energy adoption scenarios were shaped by social factors, and that "peer effects exert a more substantial influence on adoption decisions than social norms," pointing to direct interpersonal influence rather than normative pressure. Mochi et al. (2025) reinforce these findings, showing that social comparison and group identification could lead to a decrease of 6.7\% to 11\% in total consumption. Although these studies refer to general electricity usage and appliances rather than changing usage times, it implies that electricity usage is influenced in some way through the social network.  Falandays and Smaldino (2022) provide a general theoretical background using \textit{Cultural Attractor Theory} and ABM. Through repeated social transmission, individually variable behaviors tend to converge toward stable attractor points within a population. This \textit{cognitive alignment} emerges from accumulated exposure to neighbors' patterns rather than from explicit coordination. This suggests that even modest social influence, can produce durable shifts in behavioral norms rather than short-lived nudge effects.

This relatively optimistic and homogenetic view from Falandays and Smaldino (2022) on social influence does sit in tension with the habit literature by Häseler and Wulf (2024). \cite{mochi_social_2025} observe that social nudge effects vary significantly across different populations. In collectivist societies, peer comparison effectively drives behavior change due to the importance of group identity. As a result, this mechanism is less effective in individualistic societies, where personal autonomy and financial self-interest dominate decision-making. 

\subsection{ABM History within Energy Research} \label{subs:abmhistory}

The electricity market forms a unique research challenge: it is simultaneously a technical system and a behavioral one, in which individual decisions aggregate into emergent system-level outcomes that neither the individual nor the system operator fully controls. According to Mundu et al. (2024), Agent-Based Modeling (ABM) is a fitting methodology for analyzing the full spectrum of the energy market. Their review mentions its use in areas from renewable energy integration to demand forecasting and policy intervention. Importantly, they note ABM’s unique capabilities in representing autonomous behaviors of individual agents within the system. The use of ABM in electricity research expanded as a consequence of market liberalization. Fatras et al. (2022) mention how the decentralization of market roles with prosumers, as described in Section~\ref{subs:shift}, created modeling problems that profit from agent-level reasoning. 

Fatras et al. (2022) developed an ABM framework to help large-scale industrial consumers, who directly bid for energy themselves rather than having an commercial DSO function as the middleman. The tool evaluates which dynamic pricing offering these buyers should participate in depending on their needs: the day-ahead market (where electricity is bought the day before delivery), the intraday market (real-time adjustments closer to delivery), and the frequency containment reserve market (where consumers automatically adjust consumption to keep the grid stable). The delivered model shows the practical role ABM can play in informing heterogenetic consumers, even after abstracting market dynamics.

Where Fatras et al. (2022) operate at the industrial scale with a single rational consumer type, Vasiljevska et al. (2017) move to the residential level and introduce the behavioral complexity. They studied residential demand within the context of EU smart meter rollouts. The authors researched how installing smart meters can affect behavior, as they provide consumers with newly accessible tools such as real-time pricing and home automation. Each household agent is built on four value dimensions: egoistic (financial savings), hedonic (comfort and convenience), biospheric (environmental protection), and altruistic (societal welfare). Agents are randomly assigned weights across these dimensions, which divides them into groups of archetypes. Then, agents evaluate electricity contracts based on their archetypal goals. Agents gain experience with their contracts over time and communicate that experience through a social network, which influences decisions within that network. The simulation tracks system-level outcomes such as electricity savings, CO\textsubscript{2} reductions, comfort, and social welfare over the simulation period. This study's use of behavioral heterogenetic archetypes and the influence of a social network inspired the approach of this thesis.

Williams et al. (2025) developed an ABM on New-Zealand neighborhood-level demand simulation. Rather than coupling behavior to archetypes, like Vasiljevska et al. (2017), Williams et al. (2025) built demand usage bottom-up through appliance usage. Each agent's daily usage schedule is generated by simulating wake times, sleep cycles, commuting, and work-from-home decisions. These routines were drawn from statistical distributions that were based on local survey data. When agents are active at home, they interact with appliances according to switch-on probability distributions per appliance. Furthermore, the model includes socioeconomic heterogeneity: agent incomes are drawn from distributions, which feeds into a wealth coefficient that influences appliance availability and usage patterns. The agents' decisions emerge from these behavioral rules rather than being imposed top-down, allowing appliance-level variability to aggregate into realistic neighborhood-scale load shapes. The load aggregate was validated against empirical DSO data, and approached high temporal similarity. The success of this model and the bottom-up approach through appliance usage make this study the main methodological framework for this project.

These models show how the field of electricity demand can be decomposed through the ABM methodology. What none of these models systematically investigates, however, is how different configurations of theoretically-grounded behavioral types interact under dynamic pricing to produce emergent aggregate outcomes. Fatras et al. (2022) restrict their model to a single consumer type, focusing on large-scale industrial consumers. Vasiljevska et al. (2017) introduce behavioral heterogeneity, but archetype membership is determined by randomly assigned value weights at initialization, and the resulting population composition is not directly controlled or examined as an experimental variable. Williams et al. (2025) achieve strong empirical fit, with demand variation staying within a few percent of real DSO data, but model behavioral variation through socioeconomic differences rather than distinct behavioral architectures. The ABM presented in this thesis attempts to build on these studies, propose a new model, and fill this gap. 

\section{Methods} \label{sec:Methods}

The ABM built for this thesis simulates agent's electricity demand based on generalizable Western-European data at a 15-minute granularity. Given a population setup, a network, a period and a random state, it produces reproducible data on (aggregate) load, energy pricing, agent behavioral variables, and system metrics. The model consists of five interacting components: the population of households defined by their behavioral architectures, appliance-level load profiles, a sampled daily network, a demand-responsive hour-price generator, and a shifting mechanism through which agents adjust their per-appliance usage probability on a daily basis. Refer to Appendix A (page~\pageref{app:a}) for a full simulation flowchart.

\subsection{Data Sources} \label{subs:data}
The OpenSTEF Liander 2024 dataset was used thrice in the building of this model. 
It provided 15-minute electricity demand measurements for 55 Dutch local distribution networks, EPEX day-ahead hourly prices and solar park generation measurements for five Dutch solar park stations (Liander, 2024). The demand measurements of a residential area were used to calibrate the model on, the EPEX prices functioned as the baseline for the pricing model, and the solar measurements were used to calculate the elasticity factor within that pricing model. 

A study by Yilmaz et al. (2017) delivered the hourly switch-on probabilities per appliance, which in turn was derived from a 2011 survey on UK appliance usage (Department of Energy \& Climate Change, 2014). While dated, the dataset’s high granularity and veracity make it a solid foundation for modern studies when calibrated to reflect current energy trends (Yilmaz et al. 2017; Williams et al., 2025). Despite its comprehensive list of appliances, Electric Vehicle (EV) charging schedules were not included. The EV switch-on probability curve was therefore estimated from Robinson et al. (2013) Figure 6, which was measured around the same time and place as the appliance data.

\subsection{The Heterogenous Population} \label{subs:population}
The population consists of N household agents. Each agent is assigned to one of the three behavioral groups, according to a configurable percentage split. Three continuous behavioral parameters are independently sampled per agent from group-specific Beta distributions: habit strength, price sensitivity, and social susceptibility. The Beta distribution is parameterized by a mean $\mu$ and concentration parameter $K$, where a higher $K$ produces a tighter distribution around $\mu$. Table \ref{tab:behavioral_params} presents the full parameter configuration for each group.

\begin{table}
    \caption{Behavioral parameter configuration per agent group. Each parameter is sampled from a Beta distribution parameterized by mean $\mu$ and concentration $K$.}
    \label{tab:behavioral_params}
    \centering
    \small
    \begin{tabular}{lrrrrrr}
        \toprule
                                  & \multicolumn{2}{c}{Habit}   & \multicolumn{2}{c}{Price}   & \multicolumn{2}{c}{Social}  \\
                                    \cmidrule(lr){2-3}             \cmidrule(lr){4-5}             \cmidrule(lr){6-7}
        Dominant Behavior Type    & $\mu$  & $K$                 & $\mu$  & $K$                 & $\mu$  & $K$                 \\
        \midrule
        Habit-driven              & 0.85   & 10                  & 0.10   & 5                   & 0.15   & 5                   \\
        \midrule
        Price-responsive          & 0.50   & 30                  & 0.85   & 10                  & 0.10   & 5                   \\
        \midrule
        Social-influenced         & 0.70   & 20                  & 0.10   & 5                   & 0.85   & 10                  \\
        \bottomrule
    \end{tabular}
\end{table}
Habit-driven agents are characterized with a high and concentrated habit strength and low responsiveness to both price and social signals. This profile follows from Häseler and Wulf (2024), who find that habitual consumers exhibit strong schedule persistence even when smart metering technology and dynamic tariffs are available. Despite their resistance to change, they still demonstrate a light capacity for some adaptation when exposed to repeated social reinforcement within their network, following the Cultural Attractor Theory from Falandays \& Smaldino (2022).

Price-responsive agents prioritize direct adjustment with the economic signal while retaining moderate habits. The higher spread of the habit parameter reflects flexibility around a behavioral baseline rather than strict routine adherence. This group exhibits the lowest level of social susceptibility, as they have already converged on their optimal, cost-minimizing strategy. Consequently, they remain somewhat indifferent to peer behavior because no superior alternative is available.

Social-influenced agents are shaped as initially habitual agents with openness to peer-driven behavioral change. Their habit strength is somewhat lower and more dispersed than the habit-driven group, indicating greater behavioral plasticity. Their price sensitivity mirrors that of habit-driven agents, meaning their primary driver for schedule adjustment is social observation rather than self-induced price engagement. 

\subsection{Appliance Usage} \label{subs:usage}

At the start of the simulation, each agent is assigned a fixed inventory of nine appliance categories: Dishwasher, Washing Machine, Tumble Dryer, Cooker, Oven, Grill, Hob, Television, Electronics. Per-agent appliance-specific power draw in kW, runtime in minutes, and maximum daily use count are sampled once from normal distributions and remain fixed throughout the simulation. All distribution parameters are adopted from Williams et al (2025). A flat, always-on baseline load represents permanently active devices, such as refrigerators, routers, and standby consumption. Following Williams et al. (2025), this load is sampled per agent from a normal distribution with $\mu = 0.40$ kW and $\sigma = 0.15$ kW, then spread uniformly across all daily 96 15-minute slots. In the analysis and agents' electricity costs this baseline is disregarded, as it is not adjustable through agent behavior and therefore irrelevant to this research. Its sole purpose is to produce a plausible aggregate load curve. 

EV ownership is assigned per agent with a 5\% probability. This percentage is based on Dutch electric vehicle data ownership (RVO, 2026) and Western European averages for 2023 (Eurostat, 2025). These specific periods were chosen for alignment with the 2024 OpenSTEF Liander dataset, thereby granting a more robust model validation against current Western-European trends. EV charging power and duration are drawn from normal distributions with $\mu = 3.3$ kW and $\mu = 3.1$ hours respectively, sourced from Robinson et al. (2013), Table 4.

Each appliance is assigned a 24-hour baseline probability distribution, which represents the likelihood of the appliance being switched on during any given hour. This baseline consists of normalized per-appliance hourly switch-on data from Yilmaz et al. (2017). These distributions are then transformed to a maximum of four Gaussian peaks per appliance, represented as [peak\_hour, peak\_magnitude, peak\_width], by fitting sums of Gaussians to the switch-on data. The number of Gaussian peaks per appliance is selected using the Bayesian Information Criterion (Schwarz, 1978). This criterion balances fit quality against model complexity and prevents unnecessary peaks. The peaks were then fitted with L-BFGS-B minimizer from the Scipy package (Virtanen et al., 2020).

Each simulated day, appliance use counts are determined through Bernoulli draws over the Yilmaz et al. (2017) hourly switch-on probabilities, capped at the agent's sampled maximum per appliance (Williams et al., 2025). For each scheduled activation, a start hour is sampled from the agent's current distribution, and a random quarter-hour offset is added for 15-minute resolution. Load at the agent's sampled power draw is spread across the corresponding runtime slots, with any activation extending past midnight carried into an overflow array. This overflow will then naturally spill over into the start of the next day. 

While this methodology established a baseline aggregate load with the expected dual-peak structure, the peak timing, magnitude, and width did not sufficiently align with empirical Liander residential grid data. Therefore, a calibration was applied to the fitted Gaussian peak parameters which achieves two critical goals. First, it ensures the model reflects realistic energy consumption patterns. Second, it leads to a valid correspondence with hourly EPEX day-ahead prices, which serves as the per-hour baseline in the pricing model. Two groups of peaks were adjusted as follows:
 
\begin{itemize}
    \item \textbf{Morning and midday peaks (8:00 to 13:00):} These peaks were shifted 2.8 hours earlier to better capture the morning demand peak at 7. The peak magnitude was increased by 0.1 and the width was reduced by 0.3, producing a somewhat sharper and more localized load. A minimum width of 0.1 was set to prevent negative distributions.
 
    \item \textbf{Evening peaks (18:00 to 22:00):} These peaks were shifted 2 hours earlier, moving them toward the late-afternoon transition period observed in the empirical data. The height was increased by 0.5 and the width by 0.5, better representing the sustained and higher-intensity consumption associated with evening residential activity.
\end{itemize}
 
The effect of this calibration on the aggregate load profile relative to the Liander reference data is shown in Figure~\ref{fig:baseline_calibration}.
\begin{figure}[h]  
    \centering
    \includegraphics[width=\textwidth]{baseline_comparison.png}
    \caption{The effect of calibration on the simulated aggregate load profile relative to the Liander reference data.}
    \label{fig:baseline_calibration}
\end{figure}


\subsection{Price Generator} \label{subs:price}

While wholesale electricity prices are influenced by complex external drivers, such as geopolitical (in)stability and industrial volatility, they remain highly correlated with demand and RES generation (Durante et al., 2022). A truly realistic EPEX day-ahead hour-price predictor extends beyond the scope of this work. This implementation will stick to the supply and demand core. 
 
Each day's aggregate simulated demand is used to estimate prices for the following day, along with a general solar-based elasticity factor. First, typical baseline prices were determined from the Liander 2024 dataset by taking the median hour price across every day of the year. Then, the baseline reference demand was established by running 50 independent 30-day simulations on a population of 500 non-shifting agents. This process yielded 1,500 hourly demand profiles, which were averaged to each agent to provide a benchmark for scale-invariant per-hour power consumption. A positive deviation between simulated and baseline demand increases the price, while a negative deviation has the opposite effect. 

These deviations are scaled according to the solar elasticity for hour $h$, reflecting hour dependent RES generation throughout the day. This factor was calculated by taking the mean generation per hour for a day across all five solar parks in the Liander 2024 dataset, then normalizing the resulting array by its maximum value so that the peak solar hour maps to 1.0 and all other hours fall proportionally below it. The elasticity is then derived through an inverted mapping of this normalized curve:

\begin{equation}
    E_{\text{solar}}(h) = E_{\text{max}} - (E_{\text{max}} - E_{\text{min}}) \times \text{Solar}_{\text{norm}}(h)
\end{equation}

\noindent where:
\begin{itemize}
    \item $E_{\text{solar}}(h)$ is the solar-dependent elasticity for hour $h$.
    \item $E_{\text{max}}$ and $E_{\text{min}}$ are the upper and lower elasticity bounds, defined as $0.6$ and $0.1$ respectively.
    \item $\text{Solar}_{\text{norm}}(h)$ is the normalized solar availability for hour $h$, such that $\text{Solar}_{\text{norm}} \in [0, 1]$.
\end{itemize}
This inversion means that hours with the highest normalized solar generation receive the lowest elasticity values, approaching the lower bound of 0.1. Conversely, low generation hours receive elasticity values close to the upper bound of 0.6. The bounds of [0.1, 0.6] are set to keep price responses within plausible ranges: a lower bound above zero prevents demand deviations from eliminating the price signal entirely during solar peak hours, while an upper bound of 0.6 prevents extreme price volatility during low-supply periods. Because the mapping is a linear transformation applied uniformly across all hours, the relative spacing of elasticity values across the day remain preserved nonetheless. 

On day 0, the historical EPEX baseline is used directly, as no prior simulated demand is available. From day 1 onward, prices are estimated per hour h using:
\begin{equation}
    P(h) = P_{base}(h) \times \left( 1 + \epsilon(h) \cdot \frac{D(h) - \bar{D}(h)}{\bar{D}(h)} \right)
\end{equation}

\noindent where:
\begin{itemize}
    \item $P(h)$ is the estimated price for hour $h$.
    \item $P_{base}(h)$ is the historical EPEX baseline price.
    \item $\epsilon(h)$ is the hourly solar elasticity constant.
    \item $D(h)$ is the mean per-agent simulated demand.
    \item $\bar{D}(h)$ is the expected neutral reference demand constant.
\end{itemize}

While this approach neglects wholesale market complexities, it provides a consistent, computationally instant, and dynamic pricing system derived from empirical baselines. 

\subsection{Social Networks} \label{subs:networks}

Each simulation run operates on a pre-built social network generated using the graph Hamiltonian approach from House (2014). A graph refers to a mathematical structure of nodes and edges, where nodes represent agents and edges represent social connections between them. A Hamiltonian is a scalar function that assigns a penalty score to any network configuration based on how far it deviates from a set of desired structural targets, such that an algorithm can navigate toward low-penalty networks. This Hamiltonian proposed by House (2014) constrains three structural properties simultaneously: 
\begin{itemize}
    \item \textbf{Link Density:} Controls the overall fraction of possible connections within the network. For a network of $N = 150$ agents with a density of $\rho = 0.08$, the expected average number of neighbors per agent is $150 \times 0.08 = 12$.
    \item \textbf{Degree Heterogeneity:} The variance-to-mean ratio of the degree distribution, determines whether agents have a similar or highly unequal numbers of connections.
    \item \textbf{Clustering Coefficient:} Captures the tendency for two neighbors of an agent to also be connected to each other, forming a triangle.
\end{itemize} 

This approach was chosen because standard alternatives typically sacrifice one structural property for another and, as stated by House (2014), they ``attribute zero likelihood to the majority of observed networks.'' This Hamiltonian avoids these limitations by constraining all three properties simultaneously, ensuring the resulting networks remain structurally realistic. This approach has been adopted in other studies involving network modeling (Reali et al., 2018; Falandays \& Smaldino, 2022), demonstrating its fit as a realistic social network construction method.

First, a link density of 0.08 is applied. While Reali et al. (2018) employed a higher link density ($\rho = 0.20$), their model incorporated language, which is a present in every interaction. This is not reflective of the network in this study, which centers on a subset of connected households who could discuss energy topics. Second, a degree heterogeneity of 1.5 is chosen to include general heterogeneity in connections. An extension to the method by House (2014) is a requirement that every agent possesses at least one connection. Finally, a clustering coefficient of 0.25 was found to be an invariant value across population sizes (Schläpfer et al., 2014).

Because evaluating all possible network configurations is computationally intractable, the method uses a Markovian Monte Carlo procedure. Starting from an empty adjacency matrix, edges are randomly proposed for addition or removal one at a time. A proposal is accepted if it reduces the Hamiltonian penalty or if a probability hits in a Bernoulli trial, giving the algorithm a chance to escape local minima. The final Hamiltionian graphs used in the simulation are pre-computed for population sizes ranging from 50 to 1,000 agents in steps of 50, with five independently seeded variants per size.

Each simulated day, a daily contact sub-network is sampled per agent from the full pregenerated network. Each agent draws a personal contact count from a normal distribution with mean 8 and standard deviation 3. This hyperparameter was informed by Mousa et al. (2021), whose contact study showed a mean of approximately 12 daily contacts per person across age groups. Their mean is downscaled in this case to reflect that electricity is a rather niche conversation topic. However, the value remains relatively high because each agent in this simulation represents an entire household rather than an individual. The selection algorithm guarantees mutual assignment, and every agent receives at least one daily contact. This daily contact network determines whose previous-day peak centers each agent observes when computing its social shift target.

\subsection{Shifting} \label{subs:shifting}

The behavioral mechanism at the core of this simulation is the daily shifting of agents' appliance distributions. Only peak centers are updated through shifting, heights and widths are fixed after initialization with the habitual shift. For the determining of habit strength, price sensitivity, and social influence parameters per heterogeneous group, refer to Table \ref{tab:behavioral_params}. 
\medskip

\textbf{Habit shift}: On day 0, all agents' peaks are adjusted once by their habit strength. Peak height is increased by $\text{habit\_str} \times \varepsilon_{\text{habit}}$ and peak width is reduced by $\text{habit\_str} \times 0.5$, with a floor of $0.2$. Highly habitual agents thus begin the simulation with taller, narrower peaks, concentrating consumption more intensely at their preferred times. 

\medskip

\textbf{Price shift}: From day 1 onward, each peak center shifts toward the nearest local price minimum detected in the current day's price curve. The price-driven shift for peak center $c$ of appliance $a$ belonging to agent $i$ is:

\begin{equation} \label{eq:price_shift}
    \Delta_{\text{price}} = \text{price\_sens}_i \times r_a \times \varepsilon_{\text{price}} \times (\hat{c}_{\text{price}} - c)
\end{equation}

\noindent where $r_a$ is a per-appliance flexibility scalar and $\hat{c}_{\text{price}}$ is the nearest local price minimum to the current center $c$. The appliance flexibility rates represent the general shifting possibility for every appliance. Washing machines and dishwashers carry the highest rate ($0.75$), as Berg et al. (2024) identify these as the most flexible appliance categories for load shifting. Cooking appliances (Oven, Cooker, Hob, and Grill) carry the lowest rate ($0.15$), reflecting the routine binding constraints of  mealtimes (Berg et al., 2024). Flexibility of EVs is set at $0.35$, as flexible EV charging options are emerging but not yet widespread (Sørensen et al., 2021). Therefore, this parameter is kept conservative. Fitting data for the remaining appliances was not present, therefore they are estimations relative to the already defined flexibilities. Tumble dryer flexibility is set at $0.60$, placing it around the highest rated category but slightly punishing it for its dependence on the washing machine.  Television and electronics are both set at $0.30$, reflecting their leisure-driven usage patterns which carry moderate flexibility. They are somewhat bounded to routine, but lack the strict contraints of cooking appliances. A full overview of flexibility rates is provided in Table~\ref{tab:shift_params}.
\medskip

\textbf{Social shift}: From day 2 onward, each peak center is additionally shifted toward the index-matched mean center of that same peak across the agent's daily contacts from the previous day. The social shift is defined as:

\begin{equation} \label{eq:social_shift}
    \Delta_{\text{social}} = \text{soc\_inf}_i \times r_a \times \varepsilon_{\text{social}} \times (\bar{c}_{\text{social}} - c)
\end{equation}

\noindent where $\bar{c}_{\text{social}}$ is the mean center of peak index $i$ across yesterday's daily contacts. If no contact used an appliance on the previous day, or if a specific peak index has no contributing contact data, the social shift for that peak is zero. The final updated center, which includes both daily shift channels, is clipped to $[0.0,\ 23.0]$, constraining all centers within day limits. 

The epsilons across their respective shifting formula were set based on face validity: social agents should remain near their initial positions when surrounded by purely habitual neighbors and all shift should occur gradually across the 30-day  period. Since the price-shift channel is a direct and noiseless source of shift, its epsilon value is kept lower than that of the social-shift channel to maintain a balanced influence between the two mechanisms. Refer to Table~\ref{tab:shift_params} for the overview of epsilons.

\begin{table}
    \caption{Per-appliance flexibility rates and epsilon parameters. Flexibility rates $r_a \in [0,1]$ scale the magnitude of both price and social shifts for each appliance. The three epsilon parameters define the strength of their respective shifting mechanism.}
    \label{tab:shift_params}
    \centering
    \small
    \begin{tabular}{lll}
        \toprule
        Appliance         & Source                        & Rate $r_a$ \\
        \midrule
        Washing Machine   & Berg et al. (2024)           & 0.75 \\
        Dishwasher        & Berg et al. (2024)           & 0.75 \\
        Tumble Dryer      & Berg et al. (2024)           & 0.60 \\
        Electric Vehicle  & Sørensen et al. (2021)       & 0.35 \\
        Television        &                              & 0.30 \\
        Electronics       &                              & 0.30 \\
        Oven              & Berg et al. (2024)           & 0.15 \\
        Cooker            & Berg et al. (2024)           & 0.15 \\
        Hob               & Berg et al. (2024)           & 0.15 \\
        Grill             & Berg et al. (2024)           & 0.15 \\
        \midrule
        Parameter                         &                & Value \\
        \midrule
        $\varepsilon_{\text{habit}}$      &                & 1.0  \\
        $\varepsilon_{\text{price}}$      &                & 0.25 \\
        $\varepsilon_{\text{social}}$     &                & 0.5  \\
        \bottomrule
    \end{tabular}
\end{table}

\subsection{Metrics} \label{subs:metrics}

Metrics are computed at the end of each simulated day and stored in two data frames: one at the agent-day level and one at the system-day level.
 
At the agent level, individual flexibility is the summed absolute center displacement across all appliance peaks on a given day. this displacement can be decomposed in price and social shift. On the final simulated day, individual adjustment is computed as the mean absolute displacement of all peak centers from their day-0 initial positions. This captures the total schedule change for an agent, accumulated over the full simulation. Normalized cost divides the agent's total appliance usage cost by total appliance energy consumption in kWh. This normalizing allows for comparison across agents regardless of absolute consumption levels. Price advantage captures the difference between the population mean price and the agent's effective price per kWh, where a positive value indicates that the agent has below-average costs.

A primary system metric is the Peak-to-Average Ratio defined as:
 
\begin{equation}
\mathrm{PAR} = \frac{\max(\text{aggregate load})}{\text{mean}(\text{aggregate load})}
\label{eq:par}
\end{equation}
 
where higher values indicate greater temporal concentration of demand and higher grid stress (Makhadmeh et al., 2019). PAR is complemented with the peak hour, which provides directional context to PAR magnitudes by indicating whether peaks occur during high- or low-supply periods.
 
To capture the structural deformation of the daily load curve relative to a no-shift baseline, the Earth Mover's Distance (EMD) is computed as the Wasserstein distance between the day-29 load profile and the median load profile of a matched no-shift baseline. The value represents the minimum amount of work needed to transform one distribution into the other.

\subsection{Technical Implementation} \label{subs:technical}

\begin{table}[h]
\centering
\begin{tabular}{llll}
\toprule
Component & Details & Version & Citation \\
\midrule
Language & Python & 3.12.8 & (Rossum, 2007) \\
\addlinespace
\multirow{4}{*}{Packages} & NumPy & 1.26.4 & (Harris et al., 2020) \\
                          & Pandas & 2.2.2 & (McKinney, 2010) \\
                          & SciPy & 1.13.1 & (Virtanen et al., 2020) \\
                          & Matplotlib & 3.9.2 & (Hunter, 2007) \\
\addlinespace
Repository & GitHub Repository & --- & (Geschiere, 2026) \\
\bottomrule
\end{tabular}
\caption{Technical Implementation Details}
\label{tab:technical_implementation}
\end{table}
Instructions for running the model, analysis notebooks, and the synthetic data used in the results can all be found in the GitHub repository.

\section{Results}
To understand model behavior and outputs, the analysis will be presented in three parts. First, a qualitative analysis of the model's behavioral dynamics will be provided using a representative population to understand the system dynamics. This is followed by the statistical analysis of within-population behavioral group differences, addressing RQ1. Finally, the system-level effects of population composition are assessed, addressing RQ2. All simulations run for 30 days at 15-minute granularity. Bonferroni correction is applied throughout to control for multiple comparisons.
Throughout the analysis, set population mixes will be referred to by their label. Refer to Table~\ref{tab:populations} for these labels and the associated population.

\begin{table}[h]
    \centering
    \small
    \caption{Overview of simulated population compositions. Composition ratios are listed as Habit/Price/Social.}
    \label{tab:populations}
    \begin{tabular}{l c l p{5.5cm}} % Increased Purpose width to 6.5cm
        \toprule
        Label & H/P/S (\%) & Scope & Purpose \\
        \midrule
        Habitual          & 90/5/5   & RQ1, 2   & Near-complete habit dominance; represents minimal DR adoption. \\
        \addlinespace
        Tipping           & 70/5/25  & RQ1, 2   & Social-leaning variant; tests peer influence under habit dominance. \\
        \addlinespace
        Progressive       & 60/30/10 & QA, RQ1, 2 & Mix with visible dynamics across all three groups. \\
        \addlinespace
        Balanced          & 50/25/25 & RQ1, 2   & Equal spread of shift-dominant groups. \\
        \addlinespace
        Price-Max.        & 20/70/10 & RQ2      & Price-responsive majority; tests system dynamics of near-optimal DR. \\
        \addlinespace
        Social-Max.       & 20/10/70 & RQ2      & Social majority; tests system dynamics of social-heavy composition. \\
        \bottomrule
    \end{tabular}
\end{table}



\subsection{Qualitative Analysis} \label{subs:QA}

In Qualitative Analysis, model behavior using the Progressive population at N = 500 is examined. This composition was selected because it contains a good share of all three behavioral groups, making the shifting dynamics across both channels visible.

\begin{figure}[h]  
    \centering
    \includegraphics[width=\textwidth]{fig1_day0_vs_day29}
    \caption{Aggregate load over a 24-hour period of the Progressive population on Day 0 and Day 29.}
    \label{fig:QAaccload}
\end{figure}

Figure~\ref{fig:QAaccload} compares the aggregate load curve on day 0, where only the initial habitual shift has been applied, to the curve on day 29, after mulltiple cycles of both price and social shifting. On day 0, the load follows the expected structure, with a morning and evening maximum at impractical times. By day 29, both peaks reduced their magnitude and shifted toward earlier hours. The evening peak became more dispersed, whereas the morning peak became narrower. The IQR bands show that variance is somewhat elevated at the peaks during initialization and more evenly distributed over the entire simulation curve at day 29.

\begin{figure}[h]  
    \centering
    \includegraphics[width=\textwidth]{fig3_price_heatmap}
    \caption{Heatmap of price/kWh by hour and simulation day}
    \label{fig:QAheatmap}
\end{figure}

Figure~\ref{fig:QAheatmap} presents the evolution of hourly prices across the 30 simulated days. Reading down a column shows the price shape on a given day, reading across a row reveals how a specific hour's price evolves over time. The expensive evening hours gradually become cheaper as aggregate demand shifts away from them. The resulting midday price increase is minimal and barely visible due to the low solar elasticity factor. The fluctuations in morning prices are slightly more pronounced, due to the high elasticity factor amplifying the impact of increased demand in early hours. 

\begin{figure}  
    \centering
    \includegraphics[width=\textwidth]{fig4_cumulative_shift_contributionsup}
    \caption{Cumulative demand-shift contributions by channel for each behavioral group}
    \label{fig:QAshiftcontr}
\end{figure}

Figure~\ref{fig:QAshiftcontr} decomposes the cumulative shift contributions into the price-driven and social-driven channels for each behavioral group. The scale difference on the y-axis across the behavioral types shows the variety in shifting distance per group, with price-responsive agents reaching the highest by a wide margin. This tracks with the noiseless nature of the price-shift, they take the optimal route to a local price minimum. In contrast, social shifting does contain noise, as agents’ networks are likely influenced by non-price-responsive agents. Social-influenced agents' social contribution outgrows price-shift as they gradually adopt the price-responsive schedule observed among their network contacts within the Progressive population. This increasing social contribution is a emergent outcome of the Cultural Attractor dynamic, where socially influenced agents are gradually pulled toward the locally dominant behavioral pattern. Habit-driven agents accumulate minimal shift through either channel.

\subsection{Group Behavior} \label{subs:rq1}

This section covers the within-population differences between behavioral groups, which addresses RQ1. Three metrics are examined: mean daily flexibility (the summed absolute peak center displacement for all appliances per day per agent), normalized cost (usage price per kWh consumed per agent), and adjustment (the mean absolute displacement from day-0 peak positions, measured on the final simulated day per agent). Each metric was evaluated across the four somewhat plausible population compositions mentioned in Table~\ref{tab:populations} at N = 500 with 30 observations per group-population cell (10 seeds × 3 network variants).

Violations of Shapiro-Wilk tests and Levene's test indicated significant departures from both normality and homoscedasticity across all group-metric combinations. Consequently, non-parametric tests were applied throughout. A Kruskal-Wallis $H$-test was used as the omnibus test for each population-metric combination, followed by pairwise Mann-Whitney $U$ tests with Bonferroni correction applied across the three metrics. Epsilon squared ($\epsilon^2$) is reported as the omnibus effect size and rank-biserial correlation ($r_{rb}$) as the pairwise effect size. Table~\ref{tab:rq1_summary} presents these group-level statistics along with means and standard deviations per metric across all four populations. 

All twelve omnibus tests ($4 \text{ populations} \times 3 \text{ metrics}$) were significant after Bonferroni correction across the three metrics ($p_{\text{bonferroni}} < 0.001$, $\alpha = .05$). For flexibility and adjustment, the Kruskal-Wallis $H$-statistic was $79.121$ across all four populations, with $\epsilon^2$ values ranging from $0.874$ to $0.886$. This consistency indicates that the three behavioral groups maintain complete rank separation on these metrics regardless of the population composition they are embedded in. For normalized cost, the separation was slightly weaker, with $\epsilon^2$ values between $0.670$ and $0.805$, reflecting greater overlap between habit-driven and social-influenced agents on this metric.

\begin{table}[p]
    \centering
    \small
    \caption{Kruskal-Wallis omnibus results and group-level summary statistics per metric across all four population compositions ($N = 500$, $n = 30$ per group-population cell). Values are reported as Mean $\pm$ SD. $\varepsilon^2$ denotes epsilon squared. All $p_{\text{bon}}$ values are Bonferroni-corrected.}
    \label{tab:rq1_summary}
    \rotatebox{270}{
        \begin{tabular}{ll ccc r l r}
            \toprule
            & & \multicolumn{3}{c}{Mean $\pm$ SD} & & & \\
            \cmidrule(lr){3-5}
            Population & Metric & Habit & Price & Social & $H$ & $p_{\text{bon}}$ & $\varepsilon^2$ \\
            \midrule
            \multirow{3}{*}{Habitual}
                & Flexibility (hrs/day)   & $1.254 \pm 0.035$ & $5.525 \pm 0.240$ & $2.165 \pm 0.289$ & 79.121 & $< .001$ & 0.886 \\
                & Norm. cost (price/kWh)  & $8.772 \pm 0.033$ & $7.929 \pm 0.055$ & $8.798 \pm 0.043$ & 61.748 & $< .001$ & 0.687 \\
                & Adjustment day-29 (hrs) & $0.915 \pm 0.022$ & $3.296 \pm 0.063$ & $1.023 \pm 0.062$ & 78.080 & $< .001$ & 0.874 \\
            \midrule
            \multirow{3}{*}{Tipping}
                & Flexibility (hrs/day)   & $1.246 \pm 0.032$ & $5.529 \pm 0.292$ & $2.052 \pm 0.108$ & 79.121 & $< .001$ & 0.886 \\
                & Norm. cost (price/kWh)  & $8.777 \pm 0.031$ & $7.946 \pm 0.050$ & $8.769 \pm 0.043$ & 60.260 & $< .001$ & 0.670 \\
                & Adjustment day-29 (hrs) & $0.917 \pm 0.021$ & $3.306 \pm 0.087$ & $1.044 \pm 0.029$ & 79.121 & $< .001$ & 0.886 \\
            \midrule
            \multirow{3}{*}{Progressive}
                & Flexibility (hrs/day)   & $1.458 \pm 0.041$ & $5.874 \pm 0.130$ & $2.612 \pm 0.121$ & 79.121 & $< .001$ & 0.886 \\
                & Norm. cost (price/kWh)  & $8.645 \pm 0.037$ & $7.912 \pm 0.022$ & $8.594 \pm 0.026$ & 72.070 & $< .001$ & 0.805 \\
                & Adjustment day-29 (hrs) & $1.120 \pm 0.028$ & $3.231 \pm 0.074$ & $1.619 \pm 0.038$ & 79.121 & $< .001$ & 0.886 \\
            \midrule
            \multirow{3}{*}{Balanced}
                & Flexibility (hrs/day)   & $1.439 \pm 0.048$ & $5.845 \pm 0.126$ & $2.468 \pm 0.057$ & 79.121 & $< .001$ & 0.886 \\
                & Norm. cost (price/kWh)  & $8.663 \pm 0.035$ & $7.933 \pm 0.036$ & $8.617 \pm 0.041$ & 68.910 & $< .001$ & 0.769 \\
                & Adjustment day-29 (hrs) & $1.103 \pm 0.032$ & $3.200 \pm 0.074$ & $1.554 \pm 0.035$ & 79.121 & $< .001$ & 0.886 \\
            \bottomrule
        \end{tabular}
    }
\end{table}


Price-responsive agents exhibit approximately four times the daily flexibility of habit-driven agents across all populations, with social-influenced agents consistently occupying an intermediate position. In terms of normalized cost, price-responsive agents achieve the lowest cost per kWh. Habit-driven and social-influenced agents repeatedly show comparable cost levels, which is consistent with the lower $\varepsilon^2$ values observed for this metric. On adjustment, the ordering mirrors that of flexibility: price-responsive agents have drifted the furthest from their day-0 schedules by day 29, habit-driven agents the least, and social-influenced agents fall in between.

Pairwise Mann-Whitney $U$ tests confirmed significant differences between all group pairs on all metrics across all four populations (all $p_{\text{Bonferroni}} < 0.001$). For flexibility and adjustment, all pairwise comparisons yielded $r_{rb} = \pm 1.0$, indicating complete non-overlap in the rank distributions. For normalized cost, complete rank separation ($r_{rb} = \pm 1.0$) was observed between habit-driven and price-responsive agents and between price-responsive and social-influenced agents. The only exception was habit-driven versus social-influenced on normalized cost in the Balanced population ($r_{rb} = \pm 0.696$), where cost distributions partially overlapped despite statistical significance.

\begin{figure}
  \centering
  \includegraphics[width=\textwidth]{fig1_rq1_headline.png}
  \caption{Behavioral group differences across populations and metrics. Each panel displays one metric across the three behavioral groups within a single population. Individual runs are shown as jittered dots color-coded by group. Mean values are marked with diamonds. Vertical bars indicate the 95\% CI around the mean, these are hardly observable given the small interval.} 
  \label{fig:rq1headline}
\end{figure} 

Figure~\ref{fig:rq1headline} illustrates these distributional differences. Slightly separated distributions on flexibility and adjustment are visible, while the overlap between habit-driven and social-influenced agents on normalized cost confirms the pattern identified in the pairwise tests. Price-responsive agents consistently form an unique group across all metrics.

\begin{figure}
  \centering
  \includegraphics[width=\textwidth]{fig3_cost_comfort_scatter.png}
  \caption{Day-29 adjustment (schedule disruption, hours) versus normalized cost (cost per kWh) at the mean and day-29.} 
  \label{fig:rq1costcomfort}
\end{figure} 

Figure~\ref{fig:rq1costcomfort} shows the cost-comfort trade-off per behavioral group across all four populations. It plots day-29 adjustment against two cost measures: the mean normalized cost averaged across all 30 simulated days (left panel) and the normalized cost on day 29 only (right panel). Adjustment serves as a proxy for schedule disruption: a higher value means an agent's appliance usage has shifted further from its day-0 timing. Immediately visible is the high dependence on population type for social-influenced agents. A price-sensitive heavy population leads to a decrease in price and increase in adjustment for social agents. Price-sensitive agents produce a highly variable point cloud with outliers in low-shift populations. Midday prices in these populations face less upward pressure from other price-sensitive agents, as larger collectives drive prices higher through collective movement. Habit-driven agents land around the same price point as socially influenced agents despite them not adjusting as much. The withdrawal of the two active shifting groups from conservative demand times may have reduced prices there, benefiting the remaining habitual agents. This might seem slightly unfair to social agents, who do not receive much financial gain after their adjustment, but their shift remains beneficial from a system stability perspective. The fitted lines give a good indication of change in the cost-comfort trade-off within a behavioral group as the population type changes, with socially-influenced agents showing the steepest slope and thus greatest sensitivity to population composition.

\subsection{System Behavior} \label{subs:rq2}
This section examines how population composition affects system-level outcomes, addressing RQ2. Six populations were tested, the four plausible compositions used in Section~\ref{subs:rq1} plus two extreme configurations, being Price-Maximalist and Social-Maximalist. Refer to Table~\ref{tab:populations} for the full overview. Checks regarding robustness of the system over network sizes occurred over N = 100, 250, 500, 750, and 1,000. Six system-level metrics were evaluated: PAR, normalized cost, price advantage, per-agent flexibility, adjustment, and Earth Mover's Distance (EMD). EMD was computed as the Wasserstein distance between the day-29 load profile and the median load profile of a matched no-shift baseline, both normalized to probability distributions over the 96 quarter-hour slots. This matched-baseline design isolates the effect of behavioral shifting from composition effects, as each baseline run shares the same population, network size, and random seed but has price and social shifting disabled ($\varepsilon_{\text{price}} = \varepsilon_{\text{social}} = 0$). Each configuration was run 30 times (10 seeds $\times$ 3 network variants).

Before examining composition effects, the effect of network size on the metrics was tested as a robustness check by pooling across all populations. This was done to verify whether the findings were not dependent on $N$. Kruskal-Wallis tests confirmed that normalized cost, price advantage, EMD, per-agent flexibility, and adjustment showed no significant dependence on N, indicating that the behavioral outcomes of the model are stable across population scales. PAR was the only significant metric affected by N ($H = 524.606$, $p_{\text{Bonferroni}} < 0.001$), showing stronger peaks at smaller network sizes compared to the mean load.

\begin{figure}
  \centering
  \includegraphics[width=\textwidth]{fig2_population_N_comparison.png}
  \caption{Change of system metric across $N$ per population.} 
  \label{fig:rq2n}
\end{figure} 

Figure~\ref{fig:rq2n} supports these outcomes. Every metric except PAR maintains a stable trend as $N$ increases. Across all plots except PAR, the same hierarchy among populations is consistently observed. The Price-Maximalist curve performs the most positive on price-related metrics, accompanied with substantial behavioral and system change to achieve those gains. The Habitual and Tipping populations are the complete opposite, showing shiftless behavior and no cost benefits. The outlying PAR requires nuance, as a higher PAR does not necessarily imply a higher grid stress if this higher peak magnitude is achieved at times of high generation. 

Population composition, tested at N = 500, produced significant differences on all six metrics (all $p_{\text{bonferroni}} < .001$). Kruskal-Wallis effect sizes ranged from $\varepsilon^2 = 0.833$ for PAR to $\varepsilon^2 = 0.965$ for per-agent flexibility, indicating that population composition explains over 83\% of rank variance across every metric. For EMD, a one-way ANOVA was applied as the distributional assumptions were met, yielding $F = 733.735$ and $\eta^2 = 0.955$. Refer to Table~\ref{tab:rq2_omnibus} for the full overview.

\begin{table}
    \caption{RQ2 omnibus test results for the effect of population composition on system-level metrics at $N = 500$ (30 days per population, 6 populations).
    Kruskal-Wallis ($H$) was used for all metrics except EMD, for which a one-way ANOVA ($F$) was applied as distributional assumptions were met.
    All $p_{\text{bon}}$ values are Bonferroni-corrected across six metrics ($\alpha = .05$).
    $\varepsilon^2$ denotes epsilon squared; $\eta^2$ denotes eta squared.}
    \label{tab:rq2_omnibus}
    \centering
    \small
    \begin{tabular}{l l r l r}
        \toprule
        Metric & Test & Statistic & $p_{\text{bon}}$ & Effect size \\
        \midrule
        PAR                   & Kruskal-Wallis & $H = 149.990$ & $< .001$ & $\varepsilon^2 = 0.833$ \\
        Normalized cost       & Kruskal-Wallis & $H = 165.162$ & $< .001$ & $\varepsilon^2 = 0.920$ \\
        Price advantage       & Kruskal-Wallis & $H = 169.059$ & $< .001$ & $\varepsilon^2 = 0.943$ \\
        Per-agent flexibility & Kruskal-Wallis & $H = 172.990$ & $< .001$ & $\varepsilon^2 = 0.965$ \\
        Adjustment            & Kruskal-Wallis & $H = 169.186$ & $< .001$ & $\varepsilon^2 = 0.944$ \\
        EMD                   & One-way ANOVA  & $F = 733.735$ & $< .001$ & $\eta^2 = 0.955$ \\
        \bottomrule
    \end{tabular}
\end{table}

Pairwise comparisons at $N$ = 500 confirmed the clear hierarchical structure among populations, as seen in Figure~\ref{fig:rq2n}. Price-Maximalist was separated from every other population on nearly every metric with $r_{rb} = \pm 1.0$, with the exception of participation rate, where effect sizes approached but did not reach perfect separation ($r_{rb} = \pm 0.996$--$0.998$). Social-Maximalist showed similarly strong separation on cost, flexibility, and adjustment metrics, but was indistinguishable from Habitual and Tipping populations on participation rate ($r_{rb} = 0.167$ and $0.311$ respectively, both $p_{\text{Bonferroni}} = 1.0$), indicating that these populations produce comparable levels of agent participation despite differing behavioral compositions. Among the four plausible compositions, two pairs were not significantly different on EMD after Bonferroni correction: Habitual versus Tipping, and Balanced versus Progressive. These non-significant pairs indicate that populations with similar shares of actively shifting agents produce comparable aggregate load deformation, when the balance between price-responsive and social-influenced agents differs slightly.

\begin{figure}
  \centering
  \includegraphics[width=\textwidth]{fig1_day29_load_by_population.png}
  \caption{Day-29 median aggregate load profiles for all six populations at N = 500 with a shaded IQR band.} 
  \label{fig:rq2day29}
\end{figure} 

Figure~\ref{fig:rq2day29} contains the day-29 aggregate load profiles for all six populations. Price-Maximalist displays the most pronounced shift and small raise of the evening peak, consistent with its dominant share of price-responsive agents. Social-Maximalist occupies an intermediate position between the Habitual-Tipping curves and the Progressive-Balanced curves. Interestingly, all populations converge around the same lightly shifted morning peak. This probably follows from the higher solar elasticity during morning hours, which dampens demand-driven price deviations and weakens the price signal that would otherwise motivate further shifting. Unlike the evening peak, where higher RES generation in reachable earlier hours across a longer period allow for more shifting room to those who desire to shift.
Additionally, Figure~\ref{fig:rq2day29} provides the context lacking in Figure~\ref{fig:rq2n} with regard to the PAR. The higher peak of Price-Maximalist occurs during high RES generation, and is therefore preferable to other peaks, despite this difference in magnitude.

\section{Discussion}

The main goal of this thesis was to evaluate how heterogeneous household behavior translates into system-level outcomes within a variable-priced electricity grid. This was assessed through an agent-based modeling approach in which three behavioral architectures, habit-driven, price-responsive, and socially-influenced, interact with social networks and dynamic pricing. These influences caused households to shift their daily appliance schedules, with respect to the agents' behavioral parameters. The results demonstrate that behavioral architecture is the dominant driver of within-population outcomes, and that population composition is the dominant driver of system-level outcomes.

\subsection{Interpretation of Key Findings}

The qualitative analysis showed that the model is capable of reproducing a dual-peak residential load aggregate, consistent with European empirical data (Anvari et al., 2022; Liander, 2024). The cumulative shift decomposition behaved according to each group's parametric design: price-responsive agents accumulated almost all shift through the price channel, while social-influenced agents showed a gradual increase in the social channel as they adopted schedules observed among price-responsive contacts. This emergent property is consistent with the Cultural Attractor Theory of Falandays and Smaldino (2022), where socially influenced agents are gradually drawn toward the locally dominant behavioral pattern. Consistent with Häseler and Wulf (2024), habitual agents did not shift much through either source. 

The within-population analysis (RQ1) revealed near-complete rank separation between the three behavioral groups on flexibility and adjustment, with epsilon-squared values consistently between 0.874 and 0.886 across all four plausible compositions. Two observations require attention. First, the consistency of the Kruskal-Wallis $H$ statistic across populations suggests that, within the ABM, behavioral architecture exerts an effect on schedule-related metrics that is largely independent of the surrounding behavioral types. Second, the partial cost overlap between habit-driven and social-influenced agents in the Balanced population suggests that social agents do not converge on the cheaper price-responsive schedule, even when price-driven agents are present in their network. This is likely due to the noisiness of the social shifting channel, where a habit majority can withhold a social agent from effectively pursuing a new strategy. Price sensitive agents do not have this disadvantage, as they directly move towards a local price minimum. The cost-comfort trade-off in Figure~\ref{fig:rq1costcomfort} adds context to this dynamic: socially-influenced agents remain anchored near the habitual region, not succeeding in closing the distance to the price-responsive agents' low-cost, high-adjustment corner. This is notable given that social agents do exhibit greater schedule displacement than habit-driven agents, yet this additional disruption yields no meaningful cost advantage. Habit-driven agents benefit from the price drop that arises once shifting agents move over to further, cheaper time slots.

The system-level analysis (RQ2) found that population composition drove every reported metric within the model. The two non-significant EMD pairs, Habitual versus Tipping and Balanced versus Progressive, suggest that aggregate load deformation depends on the share of actively shifting agents rather than on the precise balance between price-responsive and social-influenced households. This pattern was particularly notable in the Tipping composition, which produced little deformation despite its 25\% social majority. A plausible interpretation within the model is that the 5\% price-responsive portion is too small to establish the kind of behavioral attractor the social agents require to shift. Whether a similar dynamic would hold in real populations remains an open and empirically testable question that this model can only motivate, not answer.
 
\subsection{Limitations}
 
Several limitations qualify the scope of these findings. First, the parameter calibration of both the appliance-flexibility rates and the three epsilon shift constants is grounded in the literature where possible, but ultimately involves face-validity judgment. Second, the calibration of the baseline aggregate load curve to the Liander reference relied on manual peak adjustments, which degrades the benefits of using empirically grounded data. Third, the behavioral parameters of the Beta distributions in Table~\ref{tab:behavioral_params} are held fixed across all simulations. Varying these parameters would test the sensitivity of the population effects to the specific habit-strength and social-influence distributions chosen here. Fourth, the solar elasticity factor uses the same hourly mean curve every simulated day. Using an observed weather period to perform this simulation over, while predicting EPEX-prices on these forecasts, could add a realistic layer to the model.  Finally, the appliance set excludes home batteries, which are increasingly common in Western-European households. This appliance greatly personalizes how optimal behavior would be defined.  
 
\subsection{Contributions and Future Research}
 
This thesis contributes an open-source, agent-based modeling framework that bridges behavioral science and grid modeling. Where Williams et al. (2025) demonstrate the feasibility of high-fidelity neighborhood-scale demand simulation, the present work extends this line of inquiry by introducing heterogeneous behavioral architectures, a daily contact sub network, and a demand-responsive pricing mechanism. The findings suggest that the impact of dynamic-pricing roll outs may depend not only on the share of flexible consumers but on the surrounding behavioral composition that determines whether socially open households convert peer observation into schedule adjustment. This is a hypothesis the model can motivate and illustrate, and one that future empirical or experimental work could test against real consumption data.
 
Future work could extend this framework in several directions. Sensitivity analyses over the Beta-distribution parameters would establish how robust the population composition effects are to the specific behavioral assumptions made here. Adding real weather data and a richer appliance set would bring the simulation closer to the household environments it currently simplifies. The effect of behavioral composition remains untested against real consumption data. This presents a clear target for follow-up experimental or field research. A more elaborate study could also relax the assumption of static behavioral group membership, allowing agents to migrate between architectures as a function of accumulated price savings or repeated social exposure.

\section{Conclusion}

This thesis presented an agent-based model for studying how heterogeneous household behavior shapes residential demand within a variable-priced electricity system. The framework simulates how household schedules and electricity prices evolve over a 30-day period. It separates three types of behaviors: habits, price responses, and social influences. These behaviors and appliance-usage baselines are exposed to a social network, and a dynamic pricing system that reacts to demand fluctuations.
 
In response to RQ1, the simulation indicated that behavioral architecture formed a substantial driver of individual flexibility and adjustment within this model. Near complete rank separation was observed across all four tested population compositions. Under the chosen parameterization, price-responsive agents shifted approximately four times more per day than habit-driven agents. Furthermore, they achieved the lowest normalized cost, while socially-influenced agents occupied an intermediate position whose cost benefit varied depending surrounding behavioral mix.
 
In response to RQ2, population composition drove every system metric within the model, explaining over 83\% of rank variance in each. The Price-Maximalist composition produced the largest aggregate load deformation and the most favorable price advantage. Compositions with a smaller price-responsive share, most notably Tipping, produced minimal deformation despite the social-influenced proportions. This suggests that social agents in this model require sufficient exposure to price-responsive peers to shift in a meaningful way. Network size did not meaningfully affect any behavioral or cost-related metric apart from PAR, indicating that the composition effects observed are stable across the tested population scales.
 
Together, these simulation results suggest that the impact of variable-pricing structures may depend as much on the behavioral composition of the consumer population as on the mere presence of individually shifting households. This is a hypothesis the model can illustrate and motivate; empirical validation against real consumption data would be a natural next step. As Western-European energy systems continue their transition toward dynamic contracts and renewable energy integration, agent-based frameworks of this kind offer a insightful way to reason about the behavioral dynamics that dynamic pricing strategies will inevitably encounter.

REFERENCES BELOW ARE FROM THE UNI TEMPLATE 
\bibliography{references}

\section*{Appendix A} \label{app:a}

\begin{figure}[p] % [p] gives it a dedicated page for maximum size
    \centering
    \makebox[\textwidth][c]{%
        \includegraphics[width=1.15\textwidth]{ simulation_pipeline.png}
    }
    \caption{Schematic representation of the simulation pipeline.}
    \label{fig:simulation_pipeline}
\end{figure}

\end{document}ver
